% !TITLE 兵车行
% !AUTHOR 杜甫
% !DYNASTY 唐
% !ORDER 19

\begin{poem}
车辚辚\textsuperscript{1}，马萧萧\textsuperscript{2}，行人弓箭各在腰。
耶娘妻子走相送，尘埃不见咸阳桥\textsuperscript{3}。
牵衣顿足拦道哭，哭声直上干云霄\textsuperscript{4}。

道旁过者问行人，行人但云点行频\textsuperscript{5}。
或从十五北防河，便至四十西营田\textsuperscript{6}。
去时里正与裹头\textsuperscript{7}，归来头白还戍边。
边庭流血成海水，武皇开边意未已\textsuperscript{8}。
君不闻汉家山东二百州\textsuperscript{9}，千村万落生荆杞\textsuperscript{10}。
纵有健妇把锄犁，禾生陇亩无东西\textsuperscript{11}。
况复秦兵耐苦战，被驱不异犬与鸡。

长者虽有问，役夫敢申恨\textsuperscript{12}？
且如今年冬，未休关西卒\textsuperscript{13}。
县官急索租，租税从何出？
信知生男恶，反是生女好。
生女犹得嫁比邻，生男埋没随百草\textsuperscript{14}。
君不见青海头\textsuperscript{15}，古来白骨无人收。
新鬼烦冤旧鬼哭，天阴雨湿声啾啾\textsuperscript{16}。
\end{poem}

\begin{annotations}
\notetext{1}{辚辚：车轮滚动的声音。}
\notetext{2}{萧萧：战马的嘶鸣声。开篇接连用两个叠声词，制造了紧迫的听觉效果。}
\notetext{3}{咸阳桥：横跨渭水的桥，在长安城外，是西行的必经之路。送别扬起的尘土遮住了桥面。}
\notetext{4}{干云霄：直冲云霄。干（gān），触及、冲犯。哭声直冲到天上，极言哭号的惨烈。}
\notetext{5}{点行频：按名册征兵太频繁了。点行，按名册征调兵役。}
\notetext{6}{营田：屯田。在边境一边戍守一边种地。从十五岁到四十岁，整整二十五年都在边疆。}
\notetext{7}{里正与裹头：里正（村长）替他裹上头巾。古代男子成年行冠礼后用头巾裹头。去的时候年纪太小，要里正替他裹。}
\notetext{8}{武皇开边意未已：汉武帝开拓边疆的欲望没有止境。这里借汉武帝暗指唐玄宗。"意未已"是杜甫对朝廷穷兵黩武的批判。}
\notetext{9}{汉家山东二百州：以汉代唐，指华山以东的广大地区。二百州是泛指。}
\notetext{10}{荆杞：荆棘和枸杞，泛指野生灌木。千万村庄长满了荒草杂树——因为壮丁都被征走了，田地荒芜。}
\notetext{11}{禾生陇亩无东西：庄稼种在田垄上长势杂乱、分不清东西。妇女种田不如男人熟练。}
\notetext{12}{敢申恨：哪敢申诉心中的怨恨？役夫面对询问，欲言又止——连说都不敢说了。}
\notetext{13}{未休关西卒：关西（函谷关以西）的士兵还没有得到休整。今年冬天还在不停地征兵。}
\notetext{14}{埋没随百草：尸骨埋在草丛中，无人收葬。生儿子不如生女儿——这是对战争最沉痛的控诉。}
\notetext{15}{青海头：青海湖边，是唐朝与吐蕃常年交战的地方。}
\notetext{16}{啾啾：鬼魂的哭叫声。到阴雨天，新鬼旧鬼的哭声啾啾不绝。}
\end{annotations}

\begin{appreciation}
这是杜甫第一首为人民发声的诗，也是中国文学史上最早直接批判战争的长篇叙事诗之一。写这首诗时，杜甫大约四十岁，刚刚开始从个人的理想转向对社会的关怀。

天宝年间，唐玄宗好大喜功，连年对吐蕃、南诏用兵。百姓被大量征发，家破人亡，但朝廷依然不肯罢休。杜甫亲眼目睹了咸阳桥上送别的惨状，写下了这首沉痛的诗。

车辚辚，马萧萧，行人弓箭各在腰——车轮滚滚，战马嘶鸣，出征的士兵腰间挂着弓箭。开篇就是声音和画面，像一部纪录片的开场。耶娘妻子走相送，尘埃不见咸阳桥——父母妻子奔跑着来送别，扬起的尘土遮住了咸阳桥。牵衣顿足拦道哭，哭声直上干云霄——拉着衣服、跺着脚、拦在道路上痛哭，哭声直冲云霄。这是全诗最惨痛的画面，也是中国文学史上最震撼的战争场景之一。

道旁过者问行人，行人但云点行频——路边的人问出征的士兵，士兵只说征兵太频繁了。或从十五北防河，便至四十西营田——有人十五岁就去北方防河，到了四十岁还在西边屯田。去时里正与裹头，归来头白还戍边——去的时候里正给他裹头巾，回来时头发都白了还要戍边。这几句用具体的数字和细节，写出了战争对普通人一生的摧残。

信知生男恶，反是生女好——现在才知道生男孩不好，反而生女孩更好。这是何等的反讽。在重男轻女的古代社会，百姓竟然盼望生女儿，因为女儿还可以嫁给邻居，儿子却要被送上战场，尸骨埋在荒草中。这种扭曲的心理，是战争最残忍的后果：它不仅夺走人的生命，还摧毁了人最基本的价值观。

新鬼烦冤旧鬼哭，天阴雨湿声啾啾——新死的鬼魂在诉冤，旧死的鬼魂在哭泣，天阴下雨的时候，啾啾的哭声到处都是。这个结尾像鬼片一样阴森，但杜甫不是在吓人，他是在告诉我们：战争的阴魂永远不会散去，它们就飘荡在每一块被鲜血浸透的土地上。
\end{appreciation}
